\chapter{Technology Stack}\label{ch:technology}
% Owner: M2 (Siddhant L Pathak). Draft 1, 3 Nov 2026; source: blueprint §12, uv.lock and
% package-lock.json. Versions are those locked at the time of writing; P6 regenerates this
% table from the lock files at feature freeze. Rows marked "planned" are added in later phases.

\section{Technologies, libraries and services}

\subsection*{Backend (Python)}
\begin{longtable}{@{}>{\raggedright\arraybackslash}p{3.6cm}l>{\raggedright\arraybackslash}p{8.4cm}@{}}
\toprule
\textbf{Technology} & \textbf{Version} & \textbf{Purpose} \\
\midrule
\endhead
Python & 3.12 & Language for the domain, services and API. \\
uv & 0.11 & Python version, dependency and virtual-environment management (lock file). \\
FastAPI & 0.162 & HTTP API, dependency injection, OpenAPI generation, WebSockets. \\
Starlette & 1.10 & ASGI toolkit under FastAPI (middleware, CORS). \\
Uvicorn & 0.54 & ASGI server. \\
Pydantic & 2.14 & Validation of every request and response at the boundary. \\
pydantic-settings & 2.16 & Typed configuration from the environment. \\
SQLAlchemy & 2.1 & ORM and query builder behind the repository interfaces. \\
psycopg & 3.3 & PostgreSQL driver. \\
Alembic & planned (P2) & Database migrations. \\
APScheduler & planned (P2) & In-process SLA monitor job. \\
argon2-cffi, PyJWT & planned (P2) & Password hashing and short-lived access tokens. \\
Google Gen AI SDK & planned (P4) & Gemini classification, summaries and embeddings. \\
google-auth & planned (P4) & Verifying Google Sign-In ID tokens and the \texttt{hd} claim. \\
\bottomrule
\end{longtable}

\subsection*{Frontend (TypeScript)}
\begin{longtable}{@{}>{\raggedright\arraybackslash}p{3.6cm}l>{\raggedright\arraybackslash}p{8.4cm}@{}}
\toprule
\textbf{Technology} & \textbf{Version} & \textbf{Purpose} \\
\midrule
\endhead
Node.js & 22 LTS & Build and test runtime. \\
TypeScript & 5.9 & Static typing in strict mode. \\
React & 19.3 & Component-based UI for the three portals. \\
Vite & 8.3 & Development server and production build. \\
Tailwind CSS & 4.3 & Utility-first styling and per-portal themes. \\
React Router & 8.5 & Role-based routing and navigation. \\
TanStack Query & 5.105 & Server state: caching, loading and error states, retries. \\
openapi-typescript, openapi-fetch & 7.13, 0.17 & A typed API client generated from the
  backend's OpenAPI document. \\
Recharts & planned (P3) & Dashboard charts. \\
React Flow & planned (P4) & Service dependency health map. \\
Web Speech API & browser & Voice input for incident reports (Chrome, Edge). \\
\bottomrule
\end{longtable}

\subsection*{Data, services and infrastructure}
\begin{longtable}{@{}>{\raggedright\arraybackslash}p{3.6cm}l>{\raggedright\arraybackslash}p{8.4cm}@{}}
\toprule
\textbf{Technology} & \textbf{Version} & \textbf{Purpose} \\
\midrule
\endhead
SQLite & 3 & Local development and tests (\texttt{foreign\_keys=ON}, WAL). \\
PostgreSQL & 16 & Docker and deployment database; CI runs the tests on it too. \\
Google Gemini API & planned (P4) & AI classification, grounded resolution summaries and
  embeddings, with rule-based and TF-IDF fallbacks. \\
SMTP email & planned (P2) & Email notifications (console fallback when unset). \\
Telegram Bot API & planned (P4) & Urgent alerts to technicians and leads. \\
Google Identity Services & planned (P4) & Sign-in restricted to the campus domain. \\
Docker, Docker Compose & 29 & One-command startup of API, web (nginx) and PostgreSQL. \\
nginx & 1.29 & Serves the built frontend and proxies the API in Docker. \\
GitHub, GitHub Actions & -- & Version control, issues as the bug log, CI (lint, types, tests
  on SQLite and PostgreSQL, builds, end-to-end tests, report). \\
Hosting & to be chosen (P2.10) & Free-tier static host, container host and managed
  PostgreSQL. \\
\bottomrule
\end{longtable}

\subsection*{Quality and documentation tooling}
\begin{longtable}{@{}>{\raggedright\arraybackslash}p{3.6cm}l>{\raggedright\arraybackslash}p{8.4cm}@{}}
\toprule
\textbf{Tool} & \textbf{Version} & \textbf{Purpose} \\
\midrule
\endhead
pytest, pytest-cov & 9.2, 7.1 & Backend unit and integration tests with coverage. \\
Ruff & 0.17 & Python linting (including a ban on wall-clock calls) and formatting. \\
mypy & 2.4 & Static type checking in \texttt{--strict} mode. \\
import-linter & 2.15 & Machine-checked layer rules (the domain imports no framework). \\
ESLint, Prettier & 10.14, 3.9 & TypeScript linting and formatting. \\
Vitest, Testing Library & 5.1, 16.3 & Frontend component tests. \\
Playwright & 1.63 & End-to-end tests of the three demo workflows. \\
Locust & planned (P5) & Load test of the search performance target. \\
pre-commit, gitleaks & 4.7, 8.30 & Commit-time checks, including secret scanning. \\
tectonic, PlantUML & 0.17, 1.2026.9 & This report and its UML diagrams, built from sources. \\
\bottomrule
\end{longtable}

\section{Technology selection: the nine questions}

\subsection*{1. Why did you choose this programming language?}
Python~3.12 for the backend and TypeScript for the frontend.
\begin{itemize}
  \item Abstract base classes, dataclasses (frozen for value objects), enums and type hints
    express the object model concisely: interfaces, the user hierarchy, the ten state classes
    and the exception hierarchy map directly onto language features.
  \item Python has first-class libraries for the intelligent features: the Gemini SDK,
    text similarity and data handling for analytics and prediction.
  \item Python lacks compile-time access control and static typing. We compensate with
    name-mangled attributes behind read-only properties, frozen dataclasses, \texttt{mypy
    --strict} in CI, and import contracts that enforce the layering.
  \item TypeScript gives the frontend static types, and generating the API client from the
    backend's OpenAPI document keeps types consistent end to end.
\end{itemize}

\subsection*{2. Why did you choose this frontend technology?}
React with TypeScript and Tailwind CSS, built with Vite.
\begin{itemize}
  \item Components such as status and priority badges, SLA countdown chips, the filter bar and
    the notification bell are written once and reused across the three portals, while each
    portal keeps its own visual identity, as the requirements ask.
  \item The ecosystem provides what our features need: TanStack Query for loading and error
    states, Recharts for dashboards, React Flow for the dependency map, and support for an
    installable progressive web app.
  \item Vite gives a fast development loop and a small production build.
\end{itemize}

\subsection*{3. Why did you choose this backend technology?}
FastAPI.
\begin{itemize}
  \item Pydantic validates every request at the boundary, which covers the invalid-input
    requirements before data reaches the domain.
  \item The OpenAPI document is generated from the code; it gives interactive documentation
    for evaluators and the source for our typed client.
  \item Dependency injection supplies the authenticated user, role and ownership checks,
    repositories and the Clock to each route, keeping routers thin; WebSockets are built in.
\end{itemize}

\subsection*{4. Why did you select this database?}
SQLite for development and tests, PostgreSQL for Docker and deployment.
\begin{itemize}
  \item The data is relational with integrity constraints between incidents, users, services,
    SLA policies and audit records; reports need joins and aggregates.
  \item SQLite needs no installation, which makes a Windows-native setup and fast tests
    possible; PostgreSQL provides concurrent writers and row locking in deployment.
  \item SQLAlchemy~2 and Alembic sit on top, so the same code and migrations run on both, and
    CI runs the whole test suite against both databases.
\end{itemize}

\subsection*{5. Why did you choose this architecture?}
A layered, containerised modular monolith (the requirements' ``Choice~1'').
\begin{itemize}
  \item The layers separate concerns: API, services, a framework-free domain, and
    persistence and adapters that implement the domain's interfaces. Dependencies point
    inward, which import-linter checks on every commit.
  \item The domain, where the OOP design lives, can be tested in isolation with fakes and a
    controllable clock.
  \item Package boundaries match team ownership, so four people can work in parallel.
\end{itemize}

\subsection*{6. What alternatives did you consider?}
\begin{itemize}
  \item \textbf{Spring Boot with PostgreSQL:} compile-time typing and strong access control,
    but heavier setup, and Python's AI and data libraries suit our intelligent features better
    on a four-week timeline.
  \item \textbf{Node.js with MongoDB:} rejected because the audit trail and SLA reporting need
    relational integrity and joins.
  \item \textbf{Flutter with Firebase (mobile and cloud):} rejected because technician and
    administrator work needs wide screens; a responsive web app also works on phones.
  \item \textbf{Microservices:} rejected at this scale because of network and distributed-
    transaction overhead; our package boundaries allow a later split.
\end{itemize}

\subsection*{7. What are the advantages and limitations of your technology stack?}
\begin{itemize}
  \item \textbf{Advantages:} fast development, a typed API contract with generated
    documentation and client, a strong AI ecosystem, the same code on SQLite and PostgreSQL,
    and one-command startup.
  \item \textbf{Limitation: runtime typing.} Python checks types at run time only; mitigated
    by \texttt{mypy --strict} and Pydantic at the boundary.
  \item \textbf{Limitation: SQLite allows one writer at a time,} so it is used only for
    development and tests.
  \item \textbf{Limitation: external AI is a dependency;} mitigated by deterministic
    fallbacks, so classification and similarity work offline.
  \item \textbf{Limitation: the in-process scheduler} supports a single API instance only.
\end{itemize}

\subsection*{8. How does your technology support object-oriented design?}
Chapter~\ref{ch:oop} gives the evidence class by class. In short: abstract base classes define
the interfaces (strategies, notification channels, AI advisor, embedding provider, clock and
repositories); inheritance is used only for genuine is-a relationships (users, states,
exceptions); polymorphism lets strategies, channels, states and clocks be swapped; composition
ties an incident's SLA clock, notes, resolution and feedback to its lifetime; and nine design
patterns appear in named classes. Python's encapsulation rests on convention, name mangling and
immutability, which we verify with static checks rather than claim the language enforces.

\subsection*{9. How would your system evolve if the number of users increased significantly?}
\begin{itemize}
  \item Run several stateless API containers behind a load balancer.
  \item Use PostgreSQL with the existing indexes, read replicas for reports, and cached
    analytics.
  \item Move the SLA monitor and notification delivery to a task queue (for example Redis with
    Celery) so that instances do not run duplicate jobs.
  \item Fan out WebSocket messages across instances with Redis publish/subscribe.
  \item Store embeddings in pgvector for similarity search at scale, and keep pagination on
    every list.
\end{itemize}
